\section{Perverse cohomology of a dense eigencomplex}
\label{perv:section}

We now work over $\mathbb C$.  Our objective is to extract a perverse
object from an eigencomplex without losing its moving-leg eigenmaps or
fusion constraints.  Density makes the spectral parameter a free closed
orbit.  A frame along this orbit gives exactness of fixed-point Hecke
operators on a subcategory containing the eigencomplex and its
truncations; local constancy in the legs then supplies the moving
statement.  This is the argument of \cite[\S6]{CT}, with all the marked
points retained in the level structure.

Let $X/\mathbb C$ be a smooth projective connected curve of genus $g$,
let $D=x_1+\cdots+x_s$ be a nonempty reduced divisor, and put
$U=X\setminus |D|$.  Let $\cB$ be either $\cA$ or $\cA^+$ from
Section~\ref{found:level-stacks}, and let $\Lambda\subset T^*\cB$
be its generically nilpotent cone from Section~\ref{geom:section}.
All Hecke functors use the relative Satake normalization.

\begin{proposition}\label{perv:eigen}
Suppose $F\in\Dlc(\cB,E)$ is a coherent Hecke eigencomplex for a
continuous parameter
$\sigma:\pi_1^{\et}(U)\to\Gd(L)$ with Zariski-dense image, where
$L/\mathbb Q_\ell$ is finite.  Then every ${}^pH^j(F)$ has a natural
coherent eigenstructure for $\sigma$, with singular support in $\Lambda$.
The eigenmaps are natural in representations and compatible with the
unit, convolution, permutations, and fusion on every collision diagonal.
If $F\ne0$, some ${}^pH^j(F)$ is nonzero.  No global cohomological bound
on $F$ and no condition on the boundary monodromy are required.
\end{proposition}

\subsection{The Betti spectral action with all levels present}

We state the Betti input before using it.  At each $x_i$, in a formal
coordinate $t_i$, ordinary and enhanced flags correspond respectively to
\[
 \mathrm{Iw}_i=\{h\in G(\mathbb C[[t_i]]):h(0)\in B\},\qquad
 \mathrm{Iw}_i^0=\{h\in G(\mathbb C[[t_i]]):h(0)\in N_B\}.
\]
The first congruence subgroup is normal in each of these groups, and
both groups lie in the positive-loop maximal parahoric, hence in its
normalizer.  Thus the finite set of levels satisfies the
hypotheses of Nadler--Yun's level-structure results
\cite[\S6.2]{NY}.  Their level cotangent nilpotent cone is defined by
generic nilpotence and is exactly $\Lambda$, with the residue conditions
established in Section~\ref{geom:section}.

Write $\mathscr D_\Lambda(\cB)$ for the Betti category with singular
support in this cone.  The results of
\cite[Theorem~6.2.2, Proposition~6.3.2, \S6.3.6, and Theorem~6.3.7]{NY}
give the following two inputs.
First, for a finite set $I$ and representations
$\boldsymbol V=(V_i)_{i\in I}$, the moving transform satisfies
\begin{equation}\label{perv:leg-support}
 \SS\bigl(\Hecke_{I,\boldsymbol V}(K)\bigr)
       \subset \Lambda\times 0_{U^I},
       \qquad K\in\mathscr D_\Lambda(\cB).
\end{equation}
This holds on the full product, including collisions.  On smooth charts,
the isotropy proved in Section~\ref{geom:section} and the
product-stratification argument in the proof of
\cite[Proposition~6.3.2 and \S6.3.6]{NY} make
this a local product statement.  On each finite-type smooth chart,
choose a microlocal stratification whose conormals contain the
pulled-back cone.  The output is weakly constructible for its product with a
sufficiently small contractible analytic leg patch, and hence is the
pullback of a fixed-leg slice.
The transport is compatible with the Hecke tensor operations.

Second, choose $u\in U^{\mathrm{an}}$ and free generators
$\gamma_1,\ldots,\gamma_a$ of its topological fundamental group, where
$a=2g+s-1$.  Since $U^{\mathrm{an}}$ has the homotopy type of a bouquet
of these circles, its Betti stack of $\Gd$-local systems is
\[
 \mathscr L=[Y/\Gd],\qquad Y=\Gd^{a},
\]
with simultaneous conjugation.  The symmetric monoidal category
$\Perf(\mathscr L)$ acts on $\mathscr D_\Lambda(\cB)$.  For
$V\in\Rep_E(\Gd)$, the equivariant bundle
$\mathscr E_V=Y\times V$ acts by $\Hecke_{V,u}$; its tautological
automorphisms along the $\gamma_i$ act by Hecke transport
\cite[Proposition~6.3.3 and \S6.3.6]{NY}.
These statements use the Satake labeling fixed in
Section~\ref{found:section}.

The Betti categories in \cite[\S5.1 and \S6.2]{NY} do not impose global
boundedness or support in a finite-type substack.  The Hecke transforms
above preserve local bounded constructibility: on any finite-type output
chart a Satake kernel has a finite-dimensional bound, whose proper
correspondence meets only a quasi-compact part of the input stack.
Consequently the local product statement also applies to the lcc objects
and to the finite diagrams used below, locally on finite-type charts.

\begin{lemma}[The scalar spectral action]\label{perv:scalar}
Let $K\in\mathscr D_\Lambda(\cB)$ have a coherent Betti eigenstructure
with parameter represented by $y\in Y(E)$.  Every
$f\in E[Y]^{\Gd}$ acts, as a degree-zero endomorphism of the spectral
unit and hence of $K$, by $f(y)\operatorname{id}_K$.
\end{lemma}

\begin{proof}
We use the identification with excursion operators in
\cite[\S7.1 and Proposition~7.2.3]{NY}; the scalar calculation is
\cite[Lemma~6.2]{CT}.  Here is the calculation for the free group of
rank $2g+s-1$.  Matrix coefficients span $E[Y]$.  Exactness of
invariants for the reductive group $\Gd$ implies that every invariant
is a finite sum of functions
\begin{equation}\label{perv:excursion}
 (g_1,\ldots,g_a)\longmapsto
 \operatorname{Tr}_{W}\left(A\circ\bigotimes_{i=1}^{a}V_i(g_i)\right),
 \qquad W=\bigotimes_{i=1}^{a}V_i,\quad
 A\in\operatorname{End}_{\Gd}(W).
\end{equation}
Indeed products of matrix coefficients are images of
$\operatorname{End}_E(W)$ under this trace map.  A finite sum of these
finite-dimensional equivariant images contains a prescribed invariant;
exactness of invariants lifts it to invariant endomorphisms in the
corresponding finite sum of source spaces.

The excursion for \eqref{perv:excursion} inserts the coevaluation tensor
in $W\otimes W^*$, transports each $V_i$ around $\gamma_i$, applies
$A$, and evaluates.  The dual factor can remain at $u$ as an identity
leg.  Naturality, the tensor constraints, and the unit constraint of the
eigenstructure identify this operation with the same tensor calculation
on the parameter.  Its value is \eqref{perv:excursion} at $y$.
This proves the assertion for every invariant function.  In particular
the scalar follows from the coherent eigenmaps, not merely from
fixed-point eigenisomorphisms.
\end{proof}

\subsection{Equivariant frames and exactness}

The next algebraic lemma trivializes each $\mathscr E_V$ after inverting
an invariant function nonzero at a given dense parameter.
Its frames will be used only to prove
exactness.  They need not respect tensor products, since the eigenmaps
will be obtained by truncating the original coherent maps.

\begin{lemma}\label{perv:frame}
If the entries of $y\in Y(E)$ generate a Zariski-dense subgroup of
$\Gd=\mathrm{PGL}_{n,E}$, then for every nonzero representation $V$,
of dimension $b$, there are an invariant $f_V\in E[Y]^{\Gd}$ with
$f_V(y)\ne0$ and equivariant maps
\[
 \alpha_V:\mathcal O_Y^{\oplus b}\longrightarrow\mathscr E_V,
 \qquad
 \beta_V:\mathscr E_V\longrightarrow\mathcal O_Y^{\oplus b}
\]
such that
\begin{equation}\label{perv:adjugates}
 \beta_V\alpha_V=f_V\operatorname{id},\qquad
 \alpha_V\beta_V=f_V\operatorname{id}.
\end{equation}
\end{lemma}

\begin{proof}
This is the closed-orbit frame argument of \cite[Lemma~6.3]{CT}.
The stabilizer of $y$ is the centralizer of a dense subgroup, hence
$Z(\Gd)=1$.  Its orbit $O$ is closed: a one-parameter subgroup giving
a limit of the conjugated tuple forces every entry into its associated
parabolic.  Density forces that parabolic to be all of $\Gd$, so the
one-parameter subgroup is central.  The affine Hilbert--Mumford
closed-orbit criterion proves the claim.  Thus $O\simeq\Gd$.

A basis $v_1,\ldots,v_b$ at $y$ defines equivariant sections on $O$ by
$s_j(hy)=hv_j$.  Restriction
$E[Y]\otimes V\twoheadrightarrow E[O]\otimes V$ is surjective because
$O$ is closed and $Y$ is affine.  Reductivity in characteristic zero
makes invariants exact, so these sections extend equivariantly to $Y$.
Use the extensions as the columns of $\alpha_V$.
Every character of the semisimple group $\Gd$ is trivial, so
$\det V$ is trivial.  After choosing a volume form,
$f_V=\det\alpha_V$ is invariant and nonzero at $y$.
The adjugate defines a regular equivariant map $\beta_V$ on all of $Y$:
on the dense open set $f_V\ne0$ it equals $f_V\alpha_V^{-1}$, and
this identity verifies equivariance everywhere.  The two adjugate
identities give \eqref{perv:adjugates}.
\end{proof}

Let $\mathscr D^{\mathrm{lcc}}_\Lambda(\cB)$ be the lcc part of the
Betti nilpotent category.  It is stable under perverse truncation.
Locally this is the microlocal d\'evissage assertion that the singular
support of a bounded complex is the union of the singular supports of
its perverse cohomologies; it follows equivalently from the exact
perverse vanishing-cycle tests.  Smooth charts give the statement
without a uniform bound on the stack.

Fix a dense tuple $y$ as in Lemma~\ref{perv:frame}.  For each
isomorphism class of nonzero representations choose $f_V$ supplied by
that lemma for this tuple, and let $S$ be the multiplicative set they
generate.  Define
\[
 \mathscr D_S=\{K\in\mathscr D^{\mathrm{lcc}}_\Lambda(\cB):
             f_K:K\to K\text{ is invertible for every }f\in S\}.
\]
Here $f_K$ denotes the central endomorphism given by the spectral action.

\begin{lemma}\label{perv:exact}
The subcategory $\mathscr D_S$ is preserved by perverse truncations and
by all fixed-point Hecke functors.  Every fixed-point Hecke functor is
perverse $t$-exact on this subcategory.
\end{lemma}

\begin{proof}
For a degree-zero natural endomorphism $f$ of the identity, naturality
for ${}^p\tau_{\le r}K\to K$ and the universal property of truncation give
\[
 f_{{}^p\tau_{\le r}K}={}^p\tau_{\le r}(f_K).
\]
Indeed both sides are the unique endomorphism of the truncation whose
composite to $K$ equals $f_K$ after the truncation map.  The other
truncation has the analogous property.  Invertibility therefore passes
to both truncations, proving that $\mathscr D_S$ inherits the perverse
$t$-structure.  No exactness of an arbitrary spectral operator has been
assumed.

Applying the spectral action to \eqref{perv:adjugates} gives maps
\[
 K^{\oplus b}\xrightarrow{\alpha_{V,K}}\Hecke_{V,u}(K)
                \xrightarrow{\beta_{V,K}}K^{\oplus b}.
\]
Symmetric monoidality identifies $f_{\Hecke_{V,u}(K)}$ with
$\Hecke_{V,u}(f_K)$, so the Hecke functor preserves $\mathscr D_S$.
On this subcategory the adjugate identities make $f_V^{-1}\beta_{V,K}$
an inverse to $\alpha_{V,K}$.  Hence
\begin{equation}\label{perv:fixed-exact}
 \Hecke_{V,u}|_{\mathscr D_S}
          \simeq\operatorname{id}_{\mathscr D_S}^{\oplus\dim V},
\end{equation}
which proves exactness.  Hecke transport along a path from $u$ proves
the assertion at any allowed point.  The zero representation acts by
zero.  Compositions, including convolution at a shared point, are
therefore exact and preserve $\mathscr D_S$ as well.
\end{proof}

\subsection{Moving legs, collisions, and adic eigenmaps}

\begin{proof}[Proof of Proposition~\ref{perv:eigen}]
Let $K$ be the Betti realization of $F$, using
Section~\ref{found:betti}.  Theorem~\ref{det:field} and singular-support
comparison put $K$ in $\mathscr D^{\mathrm{lcc}}_\Lambda(\cB)$.
Its topological monodromy remains Zariski dense.  The topological
fundamental group has dense image in the profinite fundamental group;
a polynomial equation on the image, after adjoining its coefficients
to a finite extension of $L$, is closed for the adic topology.  An
equation on topological monodromy therefore holds on the entire adic
image.  The tuple $y$ associated with $K$ satisfies
Lemma~\ref{perv:frame}; use this tuple to choose $S$ and $\mathscr D_S$.
By Lemma~\ref{perv:scalar},
$f_{V,K}=f_V(y)\operatorname{id}_K$, so $K\in\mathscr D_S$.
Lemma~\ref{perv:exact} puts all of its perverse truncations and
cohomologies in this subcategory.

We first prove the bounds needed to truncate moving transforms.
For $A\in\mathscr D_S$ and $m=|I|$, the functor
\[
                 A\longmapsto\Hecke_{I,\boldsymbol V}(A)[m]
\]
is perverse $t$-exact.  On a contractible leg patch,
\eqref{perv:leg-support} identifies its unshifted value with a constant
family of fixed-point transforms.  The fiber is a composition of
fixed-point Hecke functors, using convolution where the legs coincide,
and hence has the bounds of $A$ by Lemma~\ref{perv:exact}.
The constant sheaf shifted by $m$ on the smooth parameter space is
perverse.  The product description therefore proves the asserted bounds
on smooth charts.  Applying this exact functor to a truncation triangle
and using its universal property yields natural isomorphisms
\begin{equation}\label{perv:moving-truncation}
 {}^pH^j\bigl(\Hecke_{I,\boldsymbol V}(A)[m]\bigr)
       \simeq\Hecke_{I,\boldsymbol V}({}^pH^jA)[m].
\end{equation}

The same bounds hold for successive operations with all previous leg
variables retained.  Locally in those variables the input is a constant
family whose fixed-leg fibers belong to $\mathscr D_S$.  Applying the
next Hecke operator uses \eqref{perv:fixed-exact} on these fibers and
\eqref{perv:leg-support} for its parameter dependence.  Thus the
truncation comparison holds throughout each composition diagram.
The shift is the dimension of the total retained parameter space,
counted once; successive operations at a shared leg add no parameter.

For clarity, let $\pi:I\twoheadrightarrow J$ be a surjection, let
$\Delta_\pi:U^J\to U^I$ be its collision diagonal, and put
$W_j=\bigotimes_{i\in\pi^{-1}(j)}V_i$.  The fusion convention gives
\[
 (1\times\Delta_\pi)^*\Hecke_{I,\boldsymbol V}(A)
                  \simeq\Hecke_{J,\boldsymbol W}(A).
\]
On the local product families just considered,
\begin{equation}\label{perv:collision-exact}
                  (1\times\Delta_\pi)^*[\,|J|-|I|\,]
\end{equation}
is perverse exact: it replaces a perverse constant factor on a smooth
$|I|$-dimensional base by the perverse constant factor on the
$|J|$-dimensional base.  This proves compatibility of
\eqref{perv:moving-truncation} with collision restriction, including
iterated collisions.  Local constancy only away from diagonals would
not suffice for this step.

These bounds construct the required comparisons in the adic category
itself.  Start with the adic truncation triangles of $F$ and apply the
adic Hecke functors.  Betti realization preserves perversity and is
conservative on the finite-type charts, so the bounds just proved show
that their images are truncation triangles.  Their universal property
therefore supplies \eqref{perv:moving-truncation} adically for $F$.
The same argument applies to the retained-leg compositions and the
shifted diagonal restriction \eqref{perv:collision-exact}.
In particular no algebraization of maps constructed only in the Betti
category is needed.

Write $\mathcal V_{I,\sigma}=\boxtimes_{i\in I}(V_i)_\sigma$.
Shift the original adic eigenmap by $m$, apply ${}^pH^j$, and use
\eqref{perv:moving-truncation}.  Exterior product with the lisse factor
$\mathcal V_{I,\sigma}[m]$ is perverse exact.  We obtain
\[
 \Hecke_{I,\boldsymbol V}({}^pH^jF)[m]
       \simeq({}^pH^jF\boxtimes\mathcal V_{I,\sigma})[m].
\]
Removing the common shift gives precisely the required eigenmap
\begin{equation}\label{perv:induced-eigen}
 \Hecke_{I,\boldsymbol V}({}^pH^jF)
                  \simeq{}^pH^jF\boxtimes\mathcal V_{I,\sigma},
\end{equation}
with no moving-leg shift.

Finally apply the functorial truncation comparisons to the diagrams of
the original coherent eigenstructure.  Naturality in representations,
the unit, and permutations commute with truncation.  The retained-leg
comparison proves compatibility with convolution, and
\eqref{perv:collision-exact} proves compatibility with fusion.
Every vertex over $U^I$, after the single shift $[|I|]$, lies in the
perverse heart: locally it is a fixed-point exact transform of
${}^pH^jF$ times a lisse parameter factor.  After collision to $U^J$,
the adjustment $[|J|-|I|]$ places all vertices in the new heart.
For objects $P,Q$ of such a heart,
$\pi_r\operatorname{Map}(P,Q)=\operatorname{Hom}(P,Q[-r])=0$ for
$r>0$.  Thus their mapping spaces are discrete, and the commutative
diagrams give all the coherent constraints.  The auxiliary choices of
frames played no role in constructing \eqref{perv:induced-eigen}.

Theorem~\ref{det:field} applied to these eigenmaps gives singular
support in $\Lambda$; one may also use perverse microlocal d\'evissage.
All bounds and comparisons were local on finite-type smooth charts.
If $F$ is nonzero, its bounded restriction to some such chart has
nonzero perverse cohomology.  Smooth perverse pullback identifies this
with a restriction of some ${}^pH^jF$, proving nonvanishing.
\end{proof}
